% Compiler: pdfLaTeX
\documentclass[11pt,a4paper]{article}

\usepackage[utf8]{inputenc}
\usepackage[T1]{fontenc}
\usepackage{lmodern}
\usepackage{microtype}
\usepackage{geometry}
\usepackage{enumitem}
\usepackage{titlesec}
\usepackage{fancyhdr}
\usepackage{hyperref}

% Document details
\newcommand{\researchtitle}{The illusion of explanatory depth}
\newcommand{\researchshorttitle}{Explanatory depth}
\newcommand{\researchdate}{2026-10-25}

% Layout
\geometry{left=2.7cm,right=2.7cm,top=2.5cm,bottom=2.5cm,headheight=14pt}
\hypersetup{hidelinks,pdftitle={\researchtitle}}
\urlstyle{same}
\frenchspacing
\setlength{\parindent}{0pt}
\setlength{\parskip}{0.55em}
\linespread{1.04}
\setlength{\emergencystretch}{1em}

\titleformat{\section}
    {\large\bfseries}{\thesection}{0.65em}{}
\titleformat{\subsection}
    {\normalsize\bfseries}{\thesubsection}{0.65em}{}
\titlespacing*{\section}{0pt}{1.15em}{0.35em}
\titlespacing*{\subsection}{0pt}{0.7em}{0.2em}
\setlist[itemize]{leftmargin=1.35em,itemsep=0.25em,topsep=0.2em}
\setlist[enumerate]{leftmargin=1.6em,itemsep=0.25em,topsep=0.2em}

\pagestyle{fancy}
\fancyhf{}
\fancyhead[L]{\small Random Research}
\fancyhead[R]{\small\researchshorttitle}
\fancyfoot[C]{\small\thepage}
\renewcommand{\headrulewidth}{0pt}
\renewcommand{\footrulewidth}{0pt}

\begin{document}

{\LARGE\bfseries\researchtitle\par}
\vspace{0.4em}
{\small\researchdate\par}

\section{Overview}

Knowing what something does can feel like knowing how it works. The illusion
of explanatory depth is the gap between how well we think we understand
a mechanism and how much we can actually explain. \cite{rozenblit}

A familiar object is a useful test. A zipper is easy to recognize and use.
Explaining how its parts join and separate is a different task.

\section{Findings}

Rozenblit and Keil asked people to assess their understanding of devices,
attempt explanations, and assess themselves again. Attempting an explanation
exposed gaps that the initial rating had missed. They also found that this
illusion was stronger for explanations than for several other kinds of
knowledge. \cite{rozenblit}

Keil's later review adds a point I would otherwise miss: people often work
with partial explanations and draw on other people or the situation around
them to fill the gaps. Everyday competence can depend on knowledge outside
our own memory. \cite{keil}

These results concern judgments about understanding. They do not tell us
that a person must explain every mechanism before using it.

\section{Example}

Imagine reading an introduction to a bicycle's gears. The explanation feels
clear, so you close the page. Later, somebody asks why changing gears affects
the effort required to pedal. You can name the chain and gears but cannot
connect their motion to the change in effort.

This hypothetical example separates recognition from explanation. Reopening
the page may make the account feel familiar again without revealing which
steps you could reproduce on your own.

\section{Discussion}

I find this idea uncomfortable in a useful way. A fluent explanation can
make a subject feel settled before I have worked through it. I would trust
that feeling less after reading about this effect.

But I do not think shallow knowledge is always a failure. For many everyday
purposes, knowing what a tool does and when to ask for help is enough.
The problem starts when a broad introduction becomes confidence about
details it never covered.

For these research notes, I would rather identify one part I still cannot
explain than keep adding terminology to make the account sound complete.

\section{Practical implications}

After reading, I would close the source and explain one mechanism in plain
language. At the first missing step, I would write down the question and
look it up.

That is a check on the scope of my understanding. A smooth explanation can
still be wrong, so the result needs comparison with a reliable source.
I would also decide how much detail the task requires before turning a
short introduction into an endless reading list.

\renewcommand{\refname}{References}
\begin{thebibliography}{9}
\small
\setlength{\itemsep}{0.35em}

\bibitem{rozenblit}
Rozenblit, L., and Keil, F. (2002).
\href{https://cogdevlab.yale.edu/sites/default/files/files/rozenblit%20%26%20keil%20%202002.pdf}
{The misunderstood limits of folk science: An illusion of explanatory depth.}
\textit{Cognitive Science}, 26(5), 521--562.

\bibitem{keil}
Keil, F. C. (2006).
\href{https://doi.org/10.1146/annurev.psych.57.102904.190100}
{Explanation and understanding.}
\textit{Annual Review of Psychology}, 57, 227--254.

\end{thebibliography}
\end{document}
